% Generated draft: requires mathematical and source review.
\documentclass[11pt]{article}
\usepackage[margin=25mm]{geometry}
\usepackage{fontspec}
\setmainfont{Latin Modern Roman}
\usepackage{amsmath,amssymb,mathtools,unicode-math}
\setmathfont{Latin Modern Math}
\usepackage{enumitem,needspace,xurl,hyperref}
\setlength{\parindent}{0pt}
\setlength{\parskip}{5pt}
\setlength{\emergencystretch}{3em}
\newcommand{\N}{\mathbb{N}}
\newcommand{\Z}{\mathbb{Z}}
\newcommand{\Q}{\mathbb{Q}}
\newcommand{\R}{\mathbb{R}}
\newcommand{\C}{\mathbb{C}}
\newcommand{\F}{\mathbb{F}}
\newcommand{\A}{\mathbb{A}}
\newcommand{\PP}{\mathbb{P}}
\newcommand{\E}{\mathbb{E}}
\newcommand{\eps}{\varepsilon}
\newcommand{\dd}{\,\mathrm{d}}
\begin{document}
\section*{UnsolvedMath ID: 800469}
\textbf{Source problem code:} catalogue entry
\subsection{Definitions and mathematical statement}
Let \(n\) denote the spatial dimension and write \(x=(x_1,\ldots,x_n)\in\R^n\). Consider a potential \(V:\R\times\R^n\to\R\) that is \(1\)-periodic in each variable, meaning
\[
V(s+k,x+m)=V(s,x)\qquad(s\in\R,\ x\in\R^n,\ k\in\Z,\ m\in\Z^n).
\]
The equation under consideration is
\[
\Delta u(x)=-\partial_sV(u(x),x),\qquad
\Delta=\sum_{i=1}^{n}\partial_{x_i}^2.
\]
The source identifies this as the Euler--Lagrange equation for the formal functional
\[
\int\left(\frac12|\nabla u(x)|^2-V(u(x),x)\right)\,\dd x.
\]
The source does not specify the regularity of \(V\) or \(u\), the integration domain, or the admissible variations defining a minimal solution. These are unresolved qualifications of the statement.

Write \(\mathbb T^{n+1}=\R^{n+1}/\Z^{n+1}\), with coordinates induced by \((s,x)\), and let \(\pi:\R^{n+1}\to\mathbb T^{n+1}\) be the quotient map. The graph of a lifted solution is
\[
\Gamma_u=\{(u(x),x):x\in\R^n\}.
\]
For \(\alpha\in\R^n\), let \(\mathcal F_\alpha(V)\) denote the source's hypothesis that \(\mathbb T^{n+1}\) admits a foliation by the projected graphs of minimal solutions of slope \(\alpha\). The source does not define minimality, slope, or the regularity of this foliation; consequently, \(\mathcal F_\alpha(V)\) retains these unresolved source terms. In particular, no requirement that a leaf be a single-valued graph over \(\mathbb T^n\) is imposed here.

\textbf{Question (Misha Bialy, Problem 4).} Is the following implication valid for the intended dimension and regularity class?
\[
\bigl(\forall\alpha\in\R^n,\ \mathcal F_\alpha(V)\bigr)
\quad\Longrightarrow\quad
\exists W:\R^n\to\R\ \ \forall(s,x)\in\R\times\R^n,
\quad V(s,x)=W(x).
\]
Thus the proposed conclusion is independence of the solution-value variable \(s\), called \(u\) in the source. When \(V\) is differentiable in that variable, this conclusion is equivalent to \(\partial_sV\equiv0\).

The source gives the following motivating case: if \(V\) is independent of \(s\), the equation becomes \(\Delta u=0\), and for every \(\alpha\in\R^n\), the solutions \(u(x)=\alpha\cdot x+c\), \(c\in\R\), give a foliation by minimal hyperplanes of slope \(\alpha\). The source also reports that a variational version for compactly supported potentials was approached in reference [bialy-mackay]; that different support assumption is background, not an additional hypothesis of this periodic question.
\subsection{Short English statement}
Misha Bialy asks whether a potential periodic in the solution value and every spatial coordinate must be independent of the solution value if, for every slope, the torus is foliated by graphs of minimizing solutions of the associated elliptic equation. The source leaves the precise meanings of minimality, slope, foliation, and the required regularity unspecified. The registry labels the problem open; this draft does not independently verify that status.
\subsection{Sources}
\begin{itemize}
\item Local UnsolvedMath registry, permanent ID \(800469\), exact source problem code \texttt{catalogue entry}. Supplied title: ``Billiards and Quantitative Symplectic Geometry --- Misha Bialy --- Problem 4.'' The record labels the problem \emph{open} and lists Albach, Barraud et al. as \texttt{proposed\_by}, with proposed year 2026. These are supplied registry attributes, not independently verified conclusions.
\item Albach, Barraud et al., \emph{Open Problems in Billiards and Quantitative Symplectic Geometry} (2026), Misha Bialy, Problem 4: \url{https://arxiv.org/abs/2602.12896}. The supplied provenance records access on 2026-10-16 and extraction by \texttt{arxiv\_latex}; no independent source visit is claimed here. Contributor attribution: Misha Bialy.
\item Reference label [moser]: the source attributes development of Aubry--Mather theory and KAM theory for this equation to J. Moser, together with V. Bangert under [bangert]. No bibliographic entry or separate link for [moser] is supplied.
\item Reference label [bangert]: V. Bangert, cited with J. Moser for the preceding theoretical background. No bibliographic entry or separate link is supplied.
\item Reference label [bialy-mackay]: the source reports an approach to a variational version for compactly supported potentials. No bibliographic entry, separate link, or precise theorem is supplied.
\item Source acknowledgment: the research was supported by ISF grant 974/24.
\end{itemize}
\end{document}
